\documentclass[10pt,a4paper]{amsart}
\usepackage{amsmath,amssymb,amsthm,graphicx}
\usepackage{verbatim}
\newenvironment{reference}{\comment}{\endcomment}
%\newenvironment{reference}{}{}
\newenvironment{slogan}{\comment}{\endcomment}
\newenvironment{history}{\comment}{\endcomment}

% For commutative diagrams we use Xy-pic
\usepackage[all]{xy}

% We use 2cell for 2-commutative diagrams.
\xyoption{2cell}
% Theorem environments.
%
\theoremstyle{plain}
\newtheorem{theorem}[subsection]{Theorem}
\newtheorem{proposition}[subsection]{Proposition}
\newtheorem{lemma}[subsection]{Lemma}

\theoremstyle{definition}
\newtheorem{definition}[subsection]{Definition}
\newtheorem{example}[subsection]{Example}
\newtheorem{exercise}[subsection]{Exercise}
\newtheorem{situation}[subsection]{Situation}

\theoremstyle{remark}
\newtheorem{remark}[subsection]{Remark}
\newtheorem{remarks}[subsection]{Remarks}

\numberwithin{equation}{subsection}

% Macros
%
\def\lim{\mathop{\mathrm{lim}}\nolimits}
\def\colim{\mathop{\mathrm{colim}}\nolimits}
\def\Spec{\mathop{\mathrm{Spec}}}
\def\Hom{\mathop{\mathrm{Hom}}\nolimits}
\def\Ext{\mathop{\mathrm{Ext}}\nolimits}
\def\SheafHom{\mathop{\mathcal{H}\!\mathit{om}}\nolimits}
\def\SheafExt{\mathop{\mathcal{E}\!\mathit{xt}}\nolimits}
\def\Sch{\mathit{Sch}}
\def\Mor{\mathop{\mathrm{Mor}}\nolimits}
\def\Ob{\mathop{\mathrm{Ob}}\nolimits}
\def\Sh{\mathop{\mathit{Sh}}\nolimits}
\def\NL{\mathop{N\!L}\nolimits}
\def\CH{\mathop{\mathrm{CH}}\nolimits}
\def\proetale{{pro\text{-}\acute{e}tale}}
\def\etale{{\acute{e}tale}}
\def\QCoh{\mathit{QCoh}}
\def\Ker{\mathop{\mathrm{Ker}}}
\def\Im{\mathop{\mathrm{Im}}}
\def\Coker{\mathop{\mathrm{Coker}}}
\def\Coim{\mathop{\mathrm{Coim}}}

% Boxtimes
%
\DeclareMathSymbol{\boxtimes}{\mathbin}{AMSa}{"02}

%
% Macros for moduli stacks/spaces
%
\def\QCohstack{\mathcal{QC}\!\mathit{oh}}
\def\Cohstack{\mathcal{C}\!\mathit{oh}}
\def\Spacesstack{\mathcal{S}\!\mathit{paces}}
\def\Quotfunctor{\mathrm{Quot}}
\def\Hilbfunctor{\mathrm{Hilb}}
\def\Curvesstack{\mathcal{C}\!\mathit{urves}}
\def\Polarizedstack{\mathcal{P}\!\mathit{olarized}}
\def\Complexesstack{\mathcal{C}\!\mathit{omplexes}}
% \Pic is the operator that assigns to X its picard group, usage \Pic(X)
% \Picardstack_{X/B} denotes the Picard stack of X over B
% \Picardfunctor_{X/B} denotes the Picard functor of X over B
\def\Pic{\mathop{\mathrm{Pic}}\nolimits}
\def\Picardstack{\mathcal{P}\!\mathit{ic}}
\def\Picardfunctor{\mathrm{Pic}}
\def\Deformationcategory{\mathcal{D}\!\mathit{ef}}
\usepackage{cleveref}
\crefname{originallemma}{original lemma}{original lemmas}
\Crefname{originallemma}{Original lemma}{Original lemmas}
\makeatletter
\newlabel{more-algebra-lemma-Rlim-perfect-gives-complete}{{09AW}{source}{}{sourcepointer-1}{}}
\newlabel{more-algebra-lemma-Rlim-perfect-gives-complete@cref}{{[originallemma][0][]09AW}{[1][0][]source}{}{}{}}
\makeatother
\crefname{originallemma}{original lemma}{original lemmas}
\Crefname{originallemma}{Original lemma}{Original lemmas}
\makeatletter
\newlabel{more-algebra-lemma-perfect}{{0658}{source}{}{sourcepointer-2}{}}
\newlabel{more-algebra-lemma-perfect@cref}{{[originallemma][0][]0658}{[1][0][]source}{}{}{}}
\makeatother
\crefname{originallemma}{original lemma}{original lemmas}
\Crefname{originallemma}{Original lemma}{Original lemmas}
\makeatletter
\newlabel{more-algebra-lemma-tor-strictly-pro-zero}{{0911}{source}{}{sourcepointer-3}{}}
\newlabel{more-algebra-lemma-tor-strictly-pro-zero@cref}{{[originallemma][0][]0911}{[1][0][]source}{}{}{}}
\makeatother
\crefname{originallemma}{original lemma}{original lemmas}
\Crefname{originallemma}{Original lemma}{Original lemmas}
\makeatletter
\newlabel{sites-cohomology-lemma-locally-constant-map}{{09BE}{source}{}{sourcepointer-4}{}}
\newlabel{sites-cohomology-lemma-locally-constant-map@cref}{{[originallemma][0][]09BE}{[1][0][]source}{}{}{}}
\makeatother
\title{Mittag-Leffler condition for the sections of a constructible complex modulo the powers of an ideal on a weakly contractible affine scheme}
\author{The Stacks Project Authors}
\date{stacks-project:proetale.tex:local-snapshot-20261012; GNU FDL 1.2 or later}
\begin{document}
\maketitle
\noindent\textit{Editorial note.} All delivered mathematical text is the author's own, copied from the pinned Stacks Project chapter file proetale.tex; no statement or proof step was rewritten, reordered or removed.\par\medskip
\noindent\textit{Editorial note.} The counted proof is the author's complete argument: the reduction to a compatible system of perfect complexes lifted to the completion, the comparison of the pro-systems of sections, and the identification of the pro-system with an I-adic quotient system via the collapsing spectral sequence.\par\medskip
\noindent\textit{Editorial note.} The label lemma-weakly-contractible-locally-constant-ML names the counted statement, and the proof environment immediately following it in the source is delivered verbatim as its proof.\par\medskip
\noindent\textit{Editorial note.} The counted statement is carried by the author's own lemma environment in the chapter source.\par\medskip
\noindent\textit{Editorial note.} Pre-open gate: python operations/agent-authoring/tools/precard\_gate.py --source proetale.tex --label lemma-weakly-contractible-locally-constant-ML --start 6204 --end 6285 --local-id 3581 returned RESULT=OK (L3581 free, occupancy FREE, public interval clear). Candidate taken from sources/stacks-proetale/candidates-halfA.json.\par\medskip
\noindent\textit{Editorial note.} The statement and proof are copied from the pinned chapter file; the cited inputs are kept as external pointers with their canonical Stacks tags.\par\medskip
\noindent\textit{Allowed standard prerequisites (editorial).}\begin{itemize}
\item Perfectness of a complex of finite tor dimension with finite type cohomology. More on Algebra, the lemma characterizing perfect complexes as those of finite tor dimension with finite type cohomology sheaves. Supplies the perfectness of the mod I reduction E\_1 of the given complex.
\item Local constancy of morphisms between complexes of locally constant sheaves. Cohomology on Sites, the lemma that a map between complexes of locally constant sheaves of modules is locally constant. Supplies the local description of the transition maps as maps of complexes E\_\{n+1\} -> E\_n.
\item Lifting a compatible system of perfect complexes modulo powers of an ideal. More on Algebra, the lemma producing a finite complex of finite projective modules over the completion with isomorphisms modulo the powers of the ideal compatible with the transition maps. Supplies the complex M\textasciicircum{}bullet over the completion and the compatible isomorphisms M\textasciicircum{}bullet/I\textasciicircum{}n M\textasciicircum{}bullet = E\_n.
\item Vanishing of the higher Tor pro-systems for a perfect complex modulo powers of an ideal. More on Algebra, the lemma that the pro-systems Tor\_\{-p\}(H\textasciicircum{}q, A/I\textasciicircum{}n A) vanish for p > 0. Supplies the collapsing of the spectral sequence computing H\textasciicircum{}0(M\textasciicircum{}bullet/I\textasciicircum{}n M\textasciicircum{}bullet), showing that this pro-system is a quotient system and hence Mittag-Leffler.
\end{itemize}\medskip
\ifdefined\hypertarget\hypertarget{sourcepointer-1}{}\fi
\noindent\textit{Source reference.} Original lemma 09AW: sources/stacks-chow/tags.txt, source lines 11774--11774. Prerequisite pointer invoked at proof line(s) 6237; it names a standard fact of a sibling chapter that the counted argument uses, and every step of the counted proof built on it is retained verbatim in the delivered body.\par\medskip
\ifdefined\hypertarget\hypertarget{sourcepointer-2}{}\fi
\noindent\textit{Source reference.} Original lemma 0658: sources/stacks-chow/tags.txt, source lines 7892--7892. Prerequisite pointer invoked at proof line(s) 6224; it names a standard fact of a sibling chapter that the counted argument uses, and every step of the counted proof built on it is retained verbatim in the delivered body.\par\medskip
\ifdefined\hypertarget\hypertarget{sourcepointer-3}{}\fi
\noindent\textit{Source reference.} Original lemma 0911: sources/stacks-chow/tags.txt, source lines 11429--11429. Prerequisite pointer invoked at proof line(s) 6280; it names a standard fact of a sibling chapter that the counted argument uses, and every step of the counted proof built on it is retained verbatim in the delivered body.\par\medskip
\ifdefined\hypertarget\hypertarget{sourcepointer-4}{}\fi
\noindent\textit{Source reference.} Original lemma 09BE: sources/stacks-chow/tags.txt, source lines 11792--11792. Prerequisite pointer invoked at proof line(s) 6232; it names a standard fact of a sibling chapter that the counted argument uses, and every step of the counted proof built on it is retained verbatim in the delivered body.\par\medskip
\begin{lemma}
\label{lemma-weakly-contractible-locally-constant-ML}
Let $X$ be a weakly contractible affine scheme. Let $\Lambda$ be a Noetherian
ring and let $I \subset \Lambda$ be an ideal. Let $K$ be an object of
$D_{cons}(X, \Lambda)$ such that
$K \otimes_\Lambda^\mathbf{L} \underline{\Lambda/I^n}$
is isomorphic in $D(X_\proetale, \Lambda/I^n)$ to a
complex of constant sheaves of $\Lambda/I^n$-modules. Then
$$
H^0(X, K \otimes_\Lambda^\mathbf{L} \Lambda/I^n)
$$
has the Mittag-Leffler condition.
\end{lemma}

\begin{proof}
Say $K \otimes_\Lambda^\mathbf{L} \underline{\Lambda/I^n}$ is isomorphic
to $\underline{E_n}$ for some object $E_n$ of $D(\Lambda/I^n)$.
Since $K \otimes_\Lambda^\mathbf{L} \underline{\Lambda/I}$ has
finite tor dimension and has finite type cohomology sheaves
we see that $E_1$ is perfect (see
More on Algebra, Lemma \ref{more-algebra-lemma-perfect}). The transition maps
$$
K \otimes_\Lambda^\mathbf{L} \underline{\Lambda/I^{n + 1}}
\to
K \otimes_\Lambda^\mathbf{L} \underline{\Lambda/I^n}
$$
locally come from (possibly many distinct) maps of complexes
$E_{n + 1} \to E_n$ in $D(\Lambda/I^{n + 1})$ see
Cohomology on Sites, Lemma \ref{sites-cohomology-lemma-locally-constant-map}.
For each $n$ choose one such map and observe that it induces
an isomorphism
$E_{n + 1} \otimes_{\Lambda/I^{n + 1}}^\mathbf{L} \Lambda/I^n \to E_n$
in $D(\Lambda/I^n)$. By
More on Algebra, Lemma \ref{more-algebra-lemma-Rlim-perfect-gives-complete}
we can find a finite complex $M^\bullet$ of finite projective
$\Lambda^\wedge$-modules and isomorphisms $M^\bullet/I^nM^\bullet \to E_n$
in $D(\Lambda/I^n)$ compatible with the transition maps.

\medskip\noindent
Now observe that for each finite collection of indices
$n > m > k$ the triple of maps
$$
H^0(X, K \otimes_\Lambda^\mathbf{L} \Lambda/I^n)
\to
H^0(X, K \otimes_\Lambda^\mathbf{L} \Lambda/I^m)
\to
H^0(X, K \otimes_\Lambda^\mathbf{L} \Lambda/I^k)
$$
is isomorphic to
$$
H^0(X, \underline{M^\bullet/I^nM^\bullet})
\to
H^0(X, \underline{M^\bullet/I^mM^\bullet})
\to
H^0(X, \underline{M^\bullet/I^kM^\bullet})
$$
Namely, choose any isomorphism
$$
\underline{M^\bullet/I^nM^\bullet} \to
K \otimes_\Lambda^\mathbf{L} \Lambda/I^n
$$
induces similar isomorphisms module $I^m$ and $I^k$ and we see that
the assertion is true. Thus to prove the lemma it suffices to show that
the system
$H^0(X, \underline{M^\bullet/I^nM^\bullet})$ has Mittag-Leffler.
Since taking sections over $X$ is exact, it suffices to prove that
the system of $\Lambda$-modules
$$
H^0(M^\bullet/I^nM^\bullet)
$$
has Mittag-Leffler. Set $A = \Lambda^\wedge$ and consider the spectral
sequence
$$
\text{Tor}_{-p}^A(H^q(M^\bullet), A/I^nA) \Rightarrow
H^{p + q}(M^\bullet/I^nM^\bullet)
$$
By More on Algebra, Lemma \ref{more-algebra-lemma-tor-strictly-pro-zero}
the pro-systems $\{\text{Tor}_{-p}^A(H^q(M^\bullet), A/I^nA)\}$ are zero
for $p > 0$. Thus the pro-system $\{H^0(M^\bullet/I^nM^\bullet)\}$
is equal to the pro-system $\{H^0(M^\bullet)/I^nH^0(M^\bullet)\}$
and the lemma is proved.
\end{proof}

\end{document}
